\scanpage{54}
\margin{53}
- \mspage{68} - \hfill n$^{\mathrm{o}}$ 269

Un vecteur isotrope n'est autre qu'un vecteur orthogonal à soi-même. Dire qu'un sous-module $F$ est isotrope signifie que $F\cap F^{\mathrm o}\ne\{0\}$, ou encore que la restriction de $\Phi$ à $F$ est dégénérée; un sous-module \uline{non isotrope} $G$ de $E$ est donc un sous module tel que la restriction de $\Phi$ à $G$ est non dégénérée. Pour qu'un sous module $F$ de $E$ soit totalement isotrope, il faut et il suffit que l'on ait $F\subset F^{\mathrm o}$. Si $F$ est un sous-module totalement isotrope de $E$, il en est de même de tout sous-module $F'$ contenu dans $F$. La somme d'une famille de sous-modules totalement isotropes et orthogonaux deux à deux est un sous-module totalement isotrope.
L'ensemble des sous-modules totalement isotropes de $E$, ordonné par inclusion, est évidemment inductif ; il en résulte que tout sous-module totalement isotrope est contenu dans un \uline{sous-module totalement dégénéré maximal}.

PROPOSITION 1 - \uline{On suppose que} $A$ \uline{est un corps}. \uline{Soit} $F$ \uline{un sous-espace non isotrope de dimension finie de} $E$ ; \uline{alors} $E$ \uline{est somme directe de} $F$ \uline{et de} $F^{\mathrm o}$. \uline{Si} $E$ \uline{est somme directe d'une famille} $(E_i)_{i\in I}$ \uline{de sous espaces deux à deux orthogonaux}, \uline{pour que} $\Phi$ \uline{soit non dégénérée}, \uline{il faut et il suffit que chacun des} $E_i$ \uline{soit non isotrope dans ces conditions l'orthogonal de} $E_i$ \uline{est le sous-espace} $\sum_{j\ne i}E_j$.

En effet, comme la restriction $\Phi'$ de $\Phi$ à $F$ est non dégénérée par hypothèse, l'application $d_{\Phi'}$ de $F$ dans son dual $F^*$ associée à droite à $\Phi'$ est injective donc bijective puisque $F$ et $F^*$ sont deux espaces de même dimension finie. Par suite, pour tout $y\in E$,
